\subsection{A checked composite-degree witness}

\paragraph{Finding A16 (exact order fifteen, with ordinary proper factors).}
Let $k=\mathbb F_2[z]/(z^3+z+1)\simeq\mathbb F_8$, and let
\[
 C/k:\quad Y^2+X^2Y=X^5+(z+1)X^3+X,\qquad
 P=\infty,\qquad R=(z^2,z).
 \tag{CE.1}
\]
Then $C$ has genus two and $2$-rank one. The class
$M=\mathcal O_C(R-P)$ has exact order $15$, and its Frobenius twist
$M^{(1)}$ lies on the Raynaud theta divisor of $C$.
There is a pointed connected cyclic etale cover of degree $15$
after a finite constant extension, of genus $16$ and $2$-rank $7$.
Its primitive conductor-$15$ factor has dimension $8$ and $2$-rank
zero, whereas its conductor-$3$ and conductor-$5$ factors are
ordinary of dimensions $2$ and $4$. Thus the whole new factor has
positive $2$-rank $6$. Finding A15 nevertheless separates every
proposed pair of distinct colliding point sections on this base.

\begin{proof}
The reciprocal coordinate $x=1/X$, $y=Y/X^2$ gives
$y^2+y=x^3+(z+1)x+x^{-1}$. Its two poles have orders $3$ and $1$,
so the genus is two and the $2$-rank is one, by the same
Artin--Schreier calculation used in A09
\cite[Lemma~2.6]{PriesZhu2012}. The point $P$ is the unique point
above $X=\infty$, with pole orders $2$ and $5$ for $X$ and $Y$.
As in A09, $\operatorname{div}(dX/X^2)=2P$; consequently
\[
 B^1_C\simeq\mathcal O_{C^{(1)}}(P^{(1)}),\qquad
 \Theta_C=C^{(1)}-P^{(1)}.
 \tag{CE.2}
\]
This twist is retained explicitly: the coefficients of (CE.1) are
in $\mathbb F_8$, not all in $\mathbb F_2$.

Define the following polynomials over $k$:
\begin{align*}
 A(X)={}&z+(z+1)X+(z^2+z)X^2+(z^2+z+1)X^3\\
        &+(z^2+z)X^4+(z^2+1)X^6+(z^2+z+1)X^7,\\
 B(X)={}&1+(z+1)X+z^2X^2+(z^2+z)X^3\\
        &+(z^2+z+1)X^4+X^5,\qquad F=A+BY.
\end{align*}
Reduction modulo $z^3+z+1$ gives the exact identities
\begin{align}
 A^2+X^2AB+(X^5+(z+1)X^3+X)B^2&=(X+z^2)^{15},
 \tag{CE.3}\\
 F(R)=0,\qquad F(\iota R)&=z\ne0,
 \quad\iota(X,Y)=(X,Y+X^2).
 \tag{CE.4}
\end{align}
The independent symbolic checker verifies every coefficient of
(CE.3) and both evaluations. Since $X=z^2$ is unramified, the norm
identity proves that the only finite zero of $F$ is $R$, with
multiplicity $15$. There are no finite poles. The leading term
$X^5Y$ has pole order $15$ at $P$, while $A$ has pole order at
most $14$. Thus
\[
                       \operatorname{div}(F)=15(R-P).
 \tag{CE.5}
\]

The class is nonzero, since a principal divisor $R-P$ would make
$C$ rational. Order $3$ is excluded by
$h^0(3P)=h^0(2P)=2$, because $2P$ is canonical. To exclude order
$5$, a function with divisor $5(R-P)$ would lie in
\[
 H^0(C,\mathcal O(5P))=\langle1,X,X^2,Y\rangle.
\]
Its $Y$ coefficient would be nonzero; after scaling it would be
$Y+a_2X^2+a_1X+a_0$. Its norm has coefficient $1$ at $X$, whereas
the monic polynomial $(X+r)^5$, $r=X(R)=z^2$, has coefficient
$r^4=z\ne1$ there. This is a contradiction. The only remaining
divisor of $15$ is $15$, proving the exact order. Equation (CE.2)
now shows that $M^{(1)}$ is on theta, as claimed. In the normalization
of A11 one obtains the original class $M$, since
$F_C^*(M^{(1)})=M^2$ and $2^{-1}\equiv8\pmod{15}$.

The normalization of $W^{15}=F$ is etale because every valuation
of $F$ is divisible by $15$. It is geometrically connected:
being a third or fifth power would make respectively
$5(R-P)$ or $3(R-P)$ principal, contrary to the exact order.
For the parameter $t=X^2/Y$ at $P$, the unit $t^{15}F$ has residue
$1$, so $tW=1$ gives a $k$-point above $P$. Over
$\mathbb F_{4096}$ the deck group is the constant group
$\mathbb Z/15$, and Riemann--Hurwitz gives genus $16$.

The units modulo $15$ consist of the two multiplication-by-$2$
orbits
\[
                \{1,2,4,8\},\qquad\{14,13,11,7\}.
\]
The theta point and its inverse give one zero Frobenius arrow on
each orbit, so the primitive factor has $2$-rank zero by the
character calculation of A11. It has dimension $\varphi(15)=8$.
The conductor-$3$ factor is ordinary of dimension $2$, by A10.

For the conductor-$5$ factor, the same norm-coefficient argument
above shows that every nonzero order-$5$ point on $C-P$ has
$X$ coordinate $1$. But at $X=1$ equation (CE.1) is
$Y^2+Y=z+1$, and $\operatorname{Tr}_{\mathbb F_8/\mathbb F_2}(z+1)=1$.
Thus there is no such $k$-point. All four nonzero powers of
$M^3$ are $k$-rational. The Abel embedding and Frobenius twisting
preserve this rationality, so none lies on the relevant theta
support. All four character arrows are therefore bijective, and
the conductor-$5$ factor is ordinary of dimension $4$.
The character decomposition, whose idempotents have odd
denominators, now gives
$f_D=1+2+4+0=7$. Finally A15 applies to the rank-one base and its
rank-zero primitive factor. This proves the point-separation claim
without asserting that the whole new factor has rank zero.
\end{proof}

\paragraph{Computation scope.}
The retained producer \texttt{construct\_theta15.py} obtains the
coefficients by Taylor linear algebra in $\mathbb F_8$.
The separate \texttt{verify\_theta15\_symbolic.py} imports no
producer arithmetic and reduces the polynomial identity using
SymPy over $\mathbb F_2$. Both pass. The example was found in a
finite search of $5488$ nonbranch rational points on $784$
displayed rank-one models over $\mathbb F_8$; this search is not
claimed to classify geometric curves or all torsion orders.
No claim is made that this curve lacks a prime-power witness:
the example verifies the composite-cover conclusion itself.

\paragraph{A smaller prime-degree comparison.}
On the rank-one curve
\[
 C_5/\mathbb F_2:\quad
 Y^2+X^2Y=X^5+X^4+X^3+X,
 \qquad R_5=(1,0),\quad P_5=\infty,
\]
the function $F_5=Y+X+1$ has norm $(X+1)^5$, vanishes at $R_5$,
has value $1$ at its conjugate, and has unique pole of order $5$
at $P_5$. Hence $\operatorname{div}(F_5)=5(R_5-P_5)$, and the
nonzero theta class has exact order $5$. Since $2$ is primitive
modulo $5$, A03 supplies a degree-$5$ separating cover. Meanwhile
$Q_5=(0,0)$ satisfies $\operatorname{div}(X)=2(Q_5-P_5)$, so
$P_5$ and $Q_5$ have the same prime-to-$2$ abelian Kummer class.
The pole calculation and theta-characteristic calculation are
identical to those above; this comparison is a direct proved
special case, not a claim of minimality among all covers.
